% Options for packages loaded elsewhere
\PassOptionsToPackage{unicode}{hyperref}
\PassOptionsToPackage{hyphens}{url}
\documentclass[
]{article}
\usepackage{xcolor}
\usepackage{amsmath,amssymb}
\setcounter{secnumdepth}{-\maxdimen} % remove section numbering
\usepackage{iftex}
\ifPDFTeX
  \usepackage[T1]{fontenc}
  \usepackage[utf8]{inputenc}
  \usepackage{textcomp} % provide euro and other symbols
\else % if luatex or xetex
  \usepackage{unicode-math} % this also loads fontspec
  \defaultfontfeatures{Scale=MatchLowercase}
  \defaultfontfeatures[\rmfamily]{Ligatures=TeX,Scale=1}
\fi
\usepackage{lmodern}
\ifPDFTeX\else
  % xetex/luatex font selection
\fi
% Use upquote if available, for straight quotes in verbatim environments
\IfFileExists{upquote.sty}{\usepackage{upquote}}{}
\IfFileExists{microtype.sty}{% use microtype if available
  \usepackage[]{microtype}
  \UseMicrotypeSet[protrusion]{basicmath} % disable protrusion for tt fonts
}{}
\makeatletter
\@ifundefined{KOMAClassName}{% if non-KOMA class
  \IfFileExists{parskip.sty}{%
    \usepackage{parskip}
  }{% else
    \setlength{\parindent}{0pt}
    \setlength{\parskip}{6pt plus 2pt minus 1pt}}
}{% if KOMA class
  \KOMAoptions{parskip=half}}
\makeatother
\ifLuaTeX
  \usepackage{luacolor}
  \usepackage[soul]{lua-ul}
\else
  \usepackage{soul}
\fi
\setlength{\emergencystretch}{3em} % prevent overfull lines
\providecommand{\tightlist}{%
  \setlength{\itemsep}{0pt}\setlength{\parskip}{0pt}}
\usepackage{bookmark}
\IfFileExists{xurl.sty}{\usepackage{xurl}}{} % add URL line breaks if available
\urlstyle{same}
\hypersetup{
  hidelinks,
  pdfcreator={LaTeX via pandoc}}

\author{}
\date{}

\begin{document}

{
\setcounter{tocdepth}{3}
\tableofcontents
}
Paul Koop

Das Pompeji-Projekt

IRARAH antwortet

Die offene Gesellschaft und ihre Feinde

Eine Erzählung aus dem Pompeji-Projekt

\emph{\hl{Wie ein Professor in einem Franziskanerhabit verschwand}}

\emph{\hl{Die Flucht über die Theiß}}

Inhaltsverzeichnis

\hyperref[spurlos-verschwunden]{\textbf{Spurlos verschwunden 2}}

\hyperref[das-kloster-in-simbach-am-inn]{\textbf{Das Kloster in Simbach am Inn 8}}

\hyperref[eingekleidet-in-der-wahrheit]{\textbf{Eingekleidet in der Wahrheit 9}}

\hyperref[flucht-uxfcber-die-theiuxdf]{\textbf{Flucht über die Theiß 13}}

\hyperref[wiedersehen-in-budapest]{\textbf{Wiedersehen in Budapest 18}}

\section{Spurlos verschwunden}\label{spurlos-verschwunden}

Der Regen fiel in sanften, gleichmäßigen Tropfen auf die Ausgrabungsstätte im Archäologischen Park von Pompeji. Die Erde unter den Füßen der Archäologen verwandelte sich allmählich in eine zähe, schlammige Masse, während das rhythmische Plätschern des Wassers die einzigen Geräusche waren, die die Stille durchbrachen. Über den Ruinen hingen dichte, graue Wolken, so tief, dass sie die umgebenden Hügel zu verschlucken schienen. Die Welt wirkte wie in einen feuchten Schleier gehüllt, die antiken Mauern und freigelegten Relikte erschienen noch vergänglicher, als könnten sie jeden Moment wieder im Boden versinken.

Dr. Leonardo Moretti, der Leiter der Ausgrabungen, stand über eine brüchige Steinmauer gelehnt, seine grauen Augen fixierten den Fortschritt der Arbeiten. Sein wettergegerbtes Gesicht war ernst, seine Gedanken wanderten zurück in die Jahrhunderte, als diese Straßen und Gebäude noch von den Bewohnern Pompejis belebt gewesen waren. Die vergangenen Tage hatten vielversprechende Funde zutage gebracht -- Fragmente von Inschriften, gut erhaltene Haushaltsgegenstände. Doch heute lag eine eigentümliche Unruhe in der Luft, die sich nicht allein durch das Wetter erklären ließ.

Plötzlich kam ein Assistent hastig auf ihn zu. Die durchnässte Kleidung klebte an seinem schmalen Körper, Schlamm spritzte bei jedem Schritt. „Dr. Moretti, die Inschrift ist fast freigelegt. Wir brauchen Martina Rossi für die Begutachtung.``

Moretti blickte von der Mauer auf. „Wo ist sie? Sie sollte längst hier sein.``

Der Assistent zuckte mit den Schultern, ein nervöses Zucken lief über sein Gesicht. „Keiner hat sie heute gesehen. Sie war auch nicht beim Frühstück.``

Ein seltsames Gefühl kroch in Moretti hoch -- als ob sich eine unsichtbare Hand um seinen Magen legte. Martina war zuverlässig, eine Frau, die jede Verabredung ernst nahm. Dass sie einfach nicht erschien, ohne Bescheid zu geben, war ungewöhnlich. Zu ungewöhnlich, um es zu ignorieren. Er sah auf seine Uhr. Es war fast Mittag. Der Regen prasselte weiter auf den Steinboden.

„Ich gehe nach ihr sehen``, sagte er, mehr zu sich selbst als zu seinem Assistenten.

Moretti eilte zu seinem Auto, das am Rand der Ausgrabungsstätte parkte. In seiner Eile vergaß er seinen Regenschirm -- ein dummer Fehler. Als er den kurzen Weg zurücklegte, bemerkte er, dass die Tropfen von den Blättern der Bäume auf seinen Kragen plätscherten. Die Kälte kroch durch den Stoff.

Im Wagen ließ er den Motor aufheulen. Die Scheibenwischer glitten über die Windschutzscheibe, das leise Kratzen mischte sich mit dem beständigen Ticken der Uhr im Armaturenbrett. Moretti dachte an Martina, an Julia. Beide waren wie Schwestern, unzertrennlich -- auch abseits der Arbeit.

Was war geschehen? Sein Kopf war voller Gedanken, die sich in chaotischen Bahnen bewegten. Hatten sie sich verletzt? Gab es einen Unfall? Er verwarf den Gedanken. Es wäre ihm zu Ohren gekommen.

Er warf einen letzten Blick auf das Gelände, dann bog er in die schmale, regennasse Straße ein, die zu ihrer Wohnung führte.

Vor dem kleinen italienischen Wohnhaus hielt er an. Die Fassade in warmem Ocker mit bröckelndem Putz war ihm vertraut -- er war oft hier gewesen. Doch an diesem regnerischen Tag wirkte das Haus anders. Die Fensterläden klapperten leise im Wind, der Regen tropfte von den Dachkanten. Eine unsichtbare Bedrohung schien über dem Ort zu liegen.

Moretti klopfte. Keine Antwort. Er klopfte lauter -- wieder nichts. Er lauschte in die Stille, hoffte, dass jeden Moment Schritte zu hören wären, dass eine der Frauen die Tür öffnen würde. Aber es blieb still.

Er spähte durch ein gekipptes Fenster. Sein Herz begann schneller zu schlagen.

Das Innere der Wohnung war ein Chaos. Kleidung lag verstreut auf den Betten, als wäre sie hastig durchwühlt worden. Ein halbgepackter Koffer stand schief im Flur, der Deckel offen. Auf dem Küchentisch lagen Papiere und Notizen verstreut, als hätte jemand in Eile nach etwas Wichtigem gesucht. Die Szenerie hatte etwas Unwirkliches -- aber die Unordnung sprach Bände über eine plöstzliche, unvorbereitete Abreise.

Moretti zog sich vom Fenster zurück. *Das passt nicht zu Martina*, dachte er. *Sie ist immer so ordentlich.*

Er rannte zu seinem Auto, ließ den Motor aufheulen und fuhr mit durchdrehenden Reifen los. Das Wasser spritzte von den Straßenrändern auf. Er musste Hilfe holen.

Es war ein grauer, verregneter Vormittag, als Moretti durch die schweren Glastüren der Polizeistation in Neapel trat. Der feuchte Geruch der Stadt haftete an seiner Kleidung, seine nassen Schuhe quietschten auf dem Marmorboden.

„Ich möchte eine Vermisstenanzeige aufgeben``, sagte er, seine Stimme drängend und zugleich erschöpft. „Zwei meiner Kollegen -- Martina Rossi und ihre Mutter Julia Rossi -- sind seit gestern Abend verschwunden. Sie sollten heute Morgen bei den Ausgrabungen sein. Ihre Wohnung \ldots`` Er suchte nach Worten. „Es ist ein Chaos. Als ob sie in Eile waren.``

Der Polizist am Empfang musterte ihn. „Wann haben Sie die beiden zuletzt gesehen?{\kern0pt}``

„Gestern Nachmittag, bei der Arbeit. Alles war normal. Sie waren wie immer in der Nähe der neuen Ausgrabungsstelle. Aber danach habe ich nichts mehr von ihnen gehört. Sie sind einfach \ldots{} verschwunden.`` Morettis Stimme brach leicht, ein seltenes Zittern, das er nicht unterdrücken konnte.

Der Polizist notierte die Angaben. „Sie sagen, ihre Wohnung sei in Unordnung gewesen? Gab es Anzeichen von Gewalt?{\kern0pt}``

Moretti schüttelte den Kopf. „Nein, nichts dergleichen. Aber es war nicht normal -- die Koffer waren halb gepackt, Kleidung überall verteilt, als ob sie überstürzt aufbrechen wollten.``

„Wir werden die Anzeige aufnehmen und mit den Ermittlungen beginnen``, sagte der Polizist. Sein Ton war beruhigend, aber Moretti spürte, dass es nicht reichte.

Er trat einen Schritt zurück, sah sich in der Station um -- Beamte eilten an ihm vorbei, Telefone klingelten. Moretti fühlte sich fehl am Platz. Hier, in der Welt der Ordnung und Vorschriften, konnte er nur hoffen, dass die Unruhe in seiner Brust bald einer Antwort weichen würde. Aber tief in seinem Inneren wusste er: Dies war erst der Anfang.

Die Polizei nahm den Fall ernst. Zwei Ermittler -- ein älterer mit grau meliertem Haar, ein jüngerer mit entschlossener Miene -- übernahmen die Akte. „Zwei Frauen, Mutter und Tochter, seit gestern Abend vermisst``, las der Ältere. „Wohnung im Chaos, keine Anzeichen eines Kampfes.`` Sie tauschten einen bedeutungsvollen Blick, dann schwang er die Autotür auf.

Sie fuhren durch die regennassen Straßen Neapels nach Pompeji. Der Himmel war weiterhin von schweren Wolken bedeckt, der Nieselregen setzte sich auf der Windschutzscheibe ab. In dieser Gegend, in der die Schatten der Camorra allgegenwärtig waren, gab es viele Gründe, warum zwei Frauen spurlos verschwinden konnten.

Als sie das Wohnhaus erreichten, stieg der Ältere aus und zog seinen Mantel enger. Die Tür der Wohnung stand einen Spalt offen -- ein Detail, das sofort seine Aufmerksamkeit auf sich zog. „Das ist nicht gut``, murmelte er. Sie betraten vorsichtig. Die Luft war abgestanden und kalt, das Licht schwach, die Stille unnatürlich.

Kleidung lag über die Betten verstreut. Der halbgepackte Koffer im Flur, der Deckel offen, ein einsamer Schuh auf dem Boden. Auf dem Küchentisch lagen zerknitterte Notizen. Der Jüngere bückte sich und hob eine auf. „Es sieht aus, als ob sie mitten in den Vorbereitungen für eine Reise waren.``

„Das sieht nach einem hastigen Aufbruch aus``, murmelte der Ältere. „Aber keine Anzeichen eines Kampfes.`` Er öffnete die Badezimmertür -- nichts Auffälliges. Alles war normal, bis auf die Unordnung.

Der Jüngere zog die Jalousien hoch. Das trübe Tageslicht fiel in den Raum. „Vielleicht sind sie geflüchtet. Es gibt keinen Hinweis auf Gewalt.``

Zurück auf der Polizeistation herrschte gespannte Atmosphäre. Die Ermittler trugen die Erkenntnisse zusammen. Ein gewöhnlicher Vermisstenfall? Es gab zu viele offene Fragen.

„Das Ganze spielt sich in einer Gegend ab, in der die Camorra ihre Finger im Spiel hat``, bemerkte der Ältere. „Es wäre nicht das erste Mal, dass Menschen spurlos verschwinden.`` Sein Kollege nickte.

Der Fall wurde an die Staatsanwaltschaft übergeben -- und an die Direzione Distrettuale Antimafia.

Im Büro der Staatsanwaltschaft herrschte bedrückende Stille. Der Staatsanwalt, ein Mann mittleren Alters mit tiefen Falten um die Augen, durchblätterte die ersten Berichte. „Ein solcher Fall in der Nähe von Pompeji könnte mit der Camorra zusammenhängen``, murmelte er. „Wir müssen jede Spur verfolgen -- Bankverbindungen, Telefonaktivitäten, Kontakte. Kein Detail ist zu unbedeutend.``

Die Ermittler standen angespannt da. Einer der jüngeren Beamten trat vor. „Die Arbeitsstelle hat die Vermisstenmeldung bestätigt. Die Kollegen bei den Ausgrabungen sind äußerst besorgt. Ich schlage vor, wir untersuchen dort noch einmal gründlich.``

Der Staatsanwalt nickte. „Gut. Und ich möchte, dass die DDA eingebunden wird.`` Er dachte einen Moment nach. „Und halten Sie nach allem Ausschau, was mit historischen Artefakten zu tun hat. In der Gegend gibt es viele wertvolle Ausgrabungen -- die organisierte Kriminalität interessiert sich dafür.``

Unterdessen bereitete ARS, die Künstliche Intelligenz, die im Geheimen für I.R.A.R.A.H operierte, ihre nächste digitale Täuschung vor. In den Netzwerken der Polizei und der Fluggesellschaften arbeiteten ihre Algorithmen schnell und effizient. Flug- und Reiseaufzeichnungen wurden manipuliert, Buchungen storniert, Passagierlisten gefälscht. Es sah jetzt so aus, als hätten Martina und Julia niemals die Stadt verlassen. Die KI verschleierte die Spuren so gründlich, dass selbst erfahrene Ermittler in einem Dickicht aus falschen Fährten gefangen waren.

Doch die Behörden gaben nicht auf. Der Verdacht, dass die beiden Historikerinnen etwas entdeckt haben könnten, das nicht ans Tageslicht kommen sollte, war zu konkret. Bei einer erneuten Durchsuchung der Wohnung fanden sie eine Visitenkarte -- *Michael Phillips, Professor an der Gregoriana, Rom*.

Die Ermittler betrachteten die Karte nachdenklich. „Wer ist dieser Mann? Warum hatten sie seine Visitenkarte?{\kern0pt}`` Die Nähe zum Vatikan ließ sie aufhorchen. Eine Verbindung zur religiösen und akademischen Elite Roms -- alles schien plötzlich wichtig.

Ein Team wurde zusammengestellt, das nach Rom geschickt wurde.

Michael Phillips saß in seinem Büro an der Gregoriana. Die Bücherregale reichten von der Decke bis zum Boden. Das Licht seiner Schreibtischlampe warf lange Schatten auf den hölzernen Tisch. Seit Tagen spürte er, wie die Anspannung um ihn herum wuchs -- wie ein Netz, das sich langsam, aber unaufhaltsam zusammenzog. Die Ermittlungen in Neapel hatten an Fahrt aufgenommen. Martina und Julia rückten ins Zentrum der Aufmerksamkeit -- und damit auch er.

Eine verschlüsselte Nachricht von ARS war eingetroffen: *„Die Eintragungen bei der Flugsicherung wurden gelöscht. Keine weiteren Spuren.``* Es war eine Erleichterung -- aber flüchtig. Der Druck wuchs mit jeder Minute. Er wusste, dass der kleinste Fehler alles gefährden könnte.

Als die Dämmerung über Rom hereinbrach, saß Michael noch an seinem Schreibtisch. Die Geräusche der Stadt drangen durch das offene Fenster -- das Summen der Straßenlaternen, Motorengeräusche in der Ferne. Ein bedrückendes Crescendo, das seine Gedanken immer schneller wirbeln ließ. Er wusste, dass die Ermittler bald kommen würden.

Dann das Klopfen an der Tür.

Ein lautes Klopfen durchbrach die Stille. Michael hatte diesen Moment erwartet, sich vorbereitet -- aber das Ziehen in seinem Magen ließ sich nicht vertreiben. Mit einem tiefen Atemzug erhob er sich, zwang sich zur Ruhe und öffnete.

Zwei Männer in dunklen Anzügen standen vor ihm. „Dr. Phillips? Wir sind von der Polizia di Stato. Es geht um das Verschwinden von Martina Rossi und Julia Rossi. Dürfen wir hereinkommen?{\kern0pt}``

Michael nickte und trat zur Seite. Sie betraten sein Büro. Er spürte die verhaltene Spannung in ihren Bewegungen -- als ob sie jede Falte in seinem Gesicht, jede unwillkürliche Geste registrierten. Er führte sie zum runden Tisch. Der ernste Polizist nahm ihm gegenüber Platz, der andere blieb am Rand des Raumes.

„Sie kennen Martina Rossi und Julia Rossi sehr gut?{\kern0pt}``, begann der Polizist.

„Ja. Ich habe mit ihnen an verschiedenen Projekten gearbeitet -- akademisch und im Rahmen von InSim.``

Der Polizist nickte knapp. „Martinas Arbeitgeber hat eine Vermisstenanzeige erstattet. Sie wurde das letzte Mal in Ihrer Nähe gesehen. Können Sie uns sagen, was am Tag ihres Verschwindens geschah?{\kern0pt}``

Michael ließ sich einen Moment Zeit. „Wir haben uns vor ihrer Abreise nach Pompeji getroffen. Sie wollten zu einem Workshop zurück nach Italien. Alles schien normal.``

Die Männer wechselten einen schnellen Blick. „Normal? Es gab keine Anzeichen, dass etwas nicht stimmte?{\kern0pt}``

Michael schüttelte den Kopf. „Nichts, was mir aufgefallen wäre.`` Aber in seinem Inneren kämpfte er. In diesem Moment summte das unsichtbare Interface von ARS leise in seinem Ohr: *„Sie überprüfen Flugzeugaufzeichnungen. Wir haben sie gelöscht. Bleib ruhig.``*

„Und was wissen Sie über den Unfall des InSim-Mercedes in der Nähe von Pompeji? Zwei Zeugen behaupten, das Fahrzeug wurde verfolgt.``

Michael spürte Schweiß auf seiner Stirn, aber er zwang sich zur Ruhe. ARS flüsterte: *„Wir haben die Überwachungsdaten bearbeitet. Sie werden keine Beweise finden.``*

„Mir ist nur bekannt, dass sie auf dem Weg zu einer Konferenz waren. Der Unfall kam unerwartet. Ich war zu diesem Zeitpunkt in Rom.``

Der Polizist beobachtete ihn. Dann zog er eine Visitenkarte hervor -- Michaels eigene. „Diese Karte wurde bei den persönlichen Gegenständen der Frauen gefunden. Können Sie das erklären?{\kern0pt}``

„Ja. Ich habe sie beiden gegeben, falls sie mich für akademische Fragen kontaktieren wollten.``

Ein tiefes Schweigen breitete sich aus. Michael spürte die Spannung, das Warten der Ermittler. Er hielt die Nerven. ARS` Stimme blieb konstant in seinem Ohr.

Nach einer endlosen Minute stand der zweite Ermittler auf und ging zum Fenster. „Sie wissen nichts von ihrem jetzigen Aufenthaltsort?{\kern0pt}``

„Leider nicht. Ich mache mir auch Sorgen um sie.``

Er kam zurück, beugte sich über den Tisch. „Sollten Sie uns etwas verheimlichen, Dr. Phillips, werden wir es herausfinden.``

Michael lächelte dünn. „Ich verstehe. Sie können mich jederzeit kontaktieren.``

Die Männer erhoben sich und gingen. Die Tür fiel ins Schloss. Michael ließ sich auf seinen Stuhl sinken.

„Sie sind gegangen. Die Spuren sind verwischt. Martina und Julia sind sicher``, flüsterte ARS.

Michael schloss die Augen und atmete tief ein. Aber die Sorge blieb -- tief in ihm vergraben. Sie alle waren nur einen kleinen Schritt vom Auffliegen entfernt.

\section{Das Kloster in Simbach am Inn}\label{das-kloster-in-simbach-am-inn}

Das Kloster der Congregatio Jesu in Simbach am Inn war fast menschenleer. Die bevorstehende Schließung machte es zum perfekten Versteck.

In einem ehemaligen Klassenzimmer saßen Martina und Julia an einem langen Holztisch. Michaels Doppelgänger lehnte an der Wand. Auf dem Laptop erschienen die Gesichter von ARS, Michael Phillips und einem I.R.A.R.A.H-Operator.

ARS` Stimme erfüllte den Raum. „Willkommen. Die Situation erfordert äußerste Diskretion.``

Michael Phillips, zugeschaltet aus Rom, wirkte angespannt. „Professor Neumann in Kassel hat mehrere Drohbriefe erhalten. Seine öffentliche Vorlesung morgen könnte zur Eskalation führen. Wir müssen ihn in Sicherheit bringen, bevor es zu spät ist.``

Auf dem Bildschirm erschienen Sicherheitsaufnahmen von Demonstranten vor der Universität Kassel -- wütende Gesichter, rote Schilder.

„Professor Neumann ist ein angesehener Wissenschaftler``, erklärte der I.R.A.R.A.H-Operator. „Er ist bekannt für seine kritische Haltung gegenüber postmodernen Strömungen und Transhumanismus. Unsere Aufgabe ist es, ihn aus der Universität zu holen und an einem sicheren Ort unterzubringen.``

Julia fragte: „Was ist, wenn sie uns unterwegs entdecken?{\kern0pt}``

Michael antwortete: „ARS hat mehrere alternative Routen identifiziert. Falls wir nicht wie geplant zur Kapelle gelangen, gibt es unterirdische Passagen aus dem Zweiten Weltkrieg. I.R.A.R.A.H hat ein sicheres Versteck in einer Kapelle am Stadtrand von Kassel vorbereitet.``

Der Doppelgänger trat von der Wand weg. „Wir sollten den Professor nach der Rettung hierher ins Kloster bringen. Die Auflösung verzögert sich. Ein sicherer Ort -- für ein paar Tage.``

ARS nickte. „Ich werde die Überwachung verstärken.``

Michael Phillips erhob seine Stimme. „Denkt daran -- diese Rettung ist mehr als eine Fluchtaktion. Es geht um die Freiheit des Denkens. Um das Recht, die Wahrheit zu suchen.``

Das Briefing endete. Die Mission begann.

\section{Eingekleidet in der Wahrheit}\label{eingekleidet-in-der-wahrheit}

Der Morgen war kühl und neblig, als Julia, Martina und Michaels Doppelgänger auf dem Campus der Universität Kassel ankamen. Ein bleierner Dunst lag über den Gebäuden, die angespannte Atmosphäre war fast greifbar. Der Tag begann nicht wie jeder andere, und das spürten sie bis in die Knochen. Bereits in der Ferne sahen sie die ersten Protestschilder, die in den trüben Himmel gereckt wurden: *„Gegen Faschismus und Wissenschaftsfeinde!{\kern0pt}``* und *„Weg mit dem Reaktionär!{\kern0pt}``* hallte es von einer Gruppe Demonstranten, die sich vor dem Hauptgebäude versammelt hatten.

„Das wird nicht leicht``, flüsterte Julia und sah sich mit besorgter Miene um. Ihre Augen scannten die Menschenmenge, die in einem wütenden Mosaik aus Gesichtern und Schildern zusammengepfercht war.

„Wir müssen uns beeilen``, sagte Michaels Doppelgänger und straffte die Schultern, die wie ein Schild gegen die aufkommende Kälte standen. „Der Professor wartet auf uns. Je schneller wir ihn von hier wegbringen, desto geringer ist die Chance, dass sie uns erkennen.``

Martina nickte und warf einen prüfenden Blick auf die grauen Wolken, die den Himmel bedeckten, als ob sie ein ungeschriebenes Omen heraufbeschworen. „I.R.A.R.A.H hat alles vorbereitet. Lasst uns keine Zeit verlieren.``

Sie schlüpften unauffällig in eines der Nebengebäude, ihre Schritte hastig, aber kontrolliert. Das Geräusch der Proteste draußen wurde lauter, doch im Inneren des Universitätsgebäudes war es still. Der Kontrast fühlte sich surreal an.

In dem spärlich eingerichteten Büro fanden sie Dr. Tobias Neumann, der in aller Eile seine Sachen in eine abgenutzte Ledertasche packte. Er war ein Mann mittleren Alters, mit scharfen Gesichtszügen und einem Ausdruck von Erschöpfung und Entschlossenheit in den Augen, die tief in ihren Höhlen lagen wie zwei Schatten in einer dunklen Gasse.

Als das Team eintrat, sah er auf und atmete erleichtert aus. „Ich habe auf euch gewartet``, sagte er und legte die Tasche ab. „Die Lage draußen spitzt sich zu. Die Demonstranten sind heute besonders aggressiv.`` Seine Stimme war ein rauer Faden, der sich durch die angespannte Luft zog.

„Wir haben alles vorbereitet``, sagte Michaels Doppelgänger ruhig und holte einen braunen Franziskanerhabit aus seiner Tasche. „Hier ist dein Ordenshabit. Sobald du ihn anziehst, bist du offiziell Bruder Timotheus -- ein Franziskaner auf dem Weg in die USA.``

Martina reichte ihm einen gefälschten Personalausweis. „I.R.A.R.A.H hat dafür gesorgt, dass du eine neue Identität hast. Dein Ordensname und deine neue Existenz sind gesichert.`` Die Worte lagen schwer auf den Schultern des Professors, der das Gewicht seiner Situation spürte.

Dr. Neumann starrte den Habit an, bevor er sich entschlossen aufrichtete. „Das ist verrückt``, murmelte er, während er sich in die braune Kutte hüllte. „Aber ich habe keine andere Wahl.``

Kaum war er fertig, führten sie ihn durch die leeren Flure des Universitätsgebäudes nach draußen. Vor der Tür parkte ein unauffälliger Lieferwagen, doch der Weg dorthin war nicht ohne Risiko. Die Proteste vor der Universität waren lauter geworden, die Menge schien angesichts der bevorstehenden Vorlesung des Professors aufgeheizt.

„Wir müssen da durch``, sagte Julia, während sie einen prüfenden Blick auf die Menge warf. „Sie dürfen nicht merken, dass du es bist. Wir haben nur wenig Zeit.``

„Ich werde mit ihm vorausgehen``, sagte Michaels Doppelgänger entschlossen. „Sie sollen glauben, wir sind eine Gruppe Franziskaner auf Pilgerreise.``

Sie bewegten sich langsam auf die Menge zu. Die Demonstranten bemerkten sie kaum -- bis ein Ruf aus der Menge ertönte: „Das ist er! Der reaktionäre Professor!{\kern0pt}``

Die Augen der Protestierenden richteten sich auf sie. Für einen Moment schien die Situation zu eskalieren. Doch Michaels Doppelgänger schob Dr. Neumann schnell in den Lieferwagen. Mit einem dumpfen Knall schloss sich die Tür, und sie fuhren los, während wütende Rufe hinter ihnen verhallten.

Auf der Autobahn nach Frankfurt war es ruhig, aber die Anspannung war noch immer spürbar, wie ein Strick, der überdehnt wird. Der Professor lehnte sich zurück und seufzte schwer. „Früher waren Diskussionen möglich``, sagte er nachdenklich. „Man konnte anderer Meinung sein, ohne beschimpft zu werden. Heute sehe ich nur noch Wut und Ignoranz. Ich hoffe, in den USA wird das anders sein.``

„Vielleicht``, sagte Michaels Doppelgänger. „Vielleicht auch nicht. Aber dort wirst du sicherer sein -- vorerst.``

Am Frankfurter Flughafen angekommen, führte das Team den Professor durch die Sicherheitskontrollen. Jeder Schritt musste präzise geplant sein, denn ein einziger Fehler konnte ihre gesamte Operation gefährden. Julia warf einen Blick über die Schulter, als der Professor seine Papiere an die Sicherheitsbeamten übergab. Die Beamten waren müde, routiniert. Sie scannten die Ausweise, nickten, winkten durch.

Nur einer zögerte -- ein junger Beamter mit Brille, der Neumanns Ausweis länger anschaute als nötig.

„Bruder Timotheus? Sie reisen nach Chicago?{\kern0pt}``

„Eine Pilgerreise``, sagte Neumann ruhig. „Zu den Franziskanern in den USA.``

Der Beamte nickte. „Gute Reise, Bruder.``

Am Gate verabschiedeten sie sich. „Sobald du in den USA bist, bist du sicher``, sagte Julia und legte dem Professor eine Hand auf die Schulter. „I.R.A.R.A.H wird sich um alles Weitere kümmern.``

Dr. Neumann nickte, eine Spur von Dankbarkeit in seinen Augen. „Ohne eure Hilfe wäre ich verloren gewesen``, sagte er leise. „Ich schulde euch mein Leben.``

Sie sahen ihm nach, als er durch das Gate ging und der Flug nach Chicago aufgerufen wurde. Ein letztes Mal warfen sie einen Blick auf den Mann, den sie gerettet hatten, bevor er hinter den Glastüren verschwand.

In den USA angekommen, wurde der Professor von einem Mitglied der franziskanischen Gemeinschaft abgeholt und zur Franciscan University of Steubenville gebracht. Das Team begleitete ihn auf dem weiten Weg zur Hochschule, deren ruhige und friedliche Atmosphäre im Kontrast zu den chaotischen Zuständen in Deutschland stand. Die Bäume, die den Campus umgaben, wirkten wie stille Wächter, die die neuen Ankömmlinge beschützten.

„Willkommen, Bruder Timotheus``, begrüßte ihn einer der Brüder mit einem warmen Lächeln. „Ihr Ruf ist Ihnen vorausgeeilt. Wir freuen uns, Sie in unserer Gemeinschaft willkommen zu heißen.``

„Es wird mir eine Ehre sein``, erwiderte der Professor und nickte leicht. Das Gewicht seiner neuen Identität fühlte sich zugleich befreiend und bedrückend an.

Am Tag nach seiner Ankunft hielt der Professor seine Antrittsvorlesung vor einer versammelten Gruppe von Franziskanern und Studenten. Das Team saß im hinteren Teil des Saals und hörte aufmerksam zu, wie der Professor seine Worte sorgfältig wählte. Die Aula war erfüllt von einer stillen Erwartung, die Luft schien zu vibrieren, als er sprach.

„In einer Zeit, in der der Mensch seine Mündigkeit an die Technologie abgibt``, begann er, „müssen wir uns wieder auf die Werte der Aufklärung und der Rationalität besinnen. Nicht die Technik wird uns retten, sondern das kritische Denken. Wir müssen den Humanismus gegen die postmodernen Strömungen verteidigen.``

Die Worte hallten im Raum wider. Ein kollektives Nicken ging durch die Reihen, die Studenten schienen den Funken der Hoffnung zu spüren, der in dem Professor steckte.

Nach der Vorlesung trafen sich Martina, Julia und Michaels Doppelgänger mit einem I.R.A.R.A.H-Agenten in einem der Hinterzimmer der Hochschule. Der Agent war ein schmaler Mann mit ernstem Gesicht.

„Es war riskant``, sagte er, seine Stimme ruhig und kontrolliert. „Aber wir haben Informationen, dass einige der Demonstranten in Deutschland das Ziel haben, auch die Franziskanische Gemeinschaft zu infiltrieren. Es könnte sein, dass sie auf Dr. Neumann aus sind.``

„Was können wir tun?{\kern0pt}``, fragte Julia besorgt.

„Wir müssen die Sicherheit des Professors gewährleisten und die Verbindung zur Community stärken. Er ist ein Ziel geworden. Die Frage ist nicht, *ob* sie versuchen werden, ihn zu finden, sondern *wann*.``

„Das bedeutet, wir müssen die Verteidigung verstärken``, fügte Martina hinzu. „Er ist nicht nur ein Professor. Er ist ein Symbol.``

Der Agent nickte. „In den nächsten Wochen wird es entscheidend sein, unsere Kommunikation abzusichern und die Bewegungen in der Community im Auge zu behalten. Es wird Zeit brauchen, um die Wellen zu glätten, die die Ereignisse in Deutschland hierher bringen.``

In den folgenden Tagen bereiteten sie sich vor. Dr. Neumann hielt weitere Vorträge, fand langsam Frieden in seiner neuen Identität. Einmal in der Woche füllte sich die Aula mit Studenten, die gespannt seinen Ideen lauschten. Hier war die Diskussion wieder möglich.

Eines Abends, nach einem besonders ermutigenden Vortrag, saß er mit seinen neuen Brüdern am Tisch. „Ich hätte nie gedacht, dass ich einmal so lebendig sein würde``, sagte er und hob sein Glas. „Auf die Freiheit des Geistes!{\kern0pt}``

„Auf die Freiheit des Geistes!{\kern0pt}``, riefen die anderen im Einklang, und ein Gefühl von Gemeinschaft erfüllte den Raum.

Dr. Neumann fühlte sich zum ersten Mal seit Langem frei. Aber er wusste auch, dass das Team, das ihn gerettet hatte, sich neuen Gefahren stellen würde -- an der ukrainisch-rumänischen Grenze, auf der Theiß, im Dunkel der Nacht.

\section{Flucht über die Theiß}\label{flucht-uxfcber-die-theiuxdf}

Es war noch früh am Morgen, als sich das Team im Briefingraum der Franziskanischen Hochschule in Steubenville versammelte. Die ersten Sonnenstrahlen durchbrachen das Morgenlicht und tauchten den Raum in ein sanftes Gold, doch die Atmosphäre war alles andere als entspannt. Die Luft war erfüllt von einer Mischung aus Konzentration und Nervosität -- jeder spürte, dass das, was vor ihnen lag, weitreichende Konsequenzen haben könnte.

An der Wand hing eine große Landkarte. Die geplante Route durch Rumänien und entlang der Theiß war in knallroter Farbe eingezeichnet. Die Kartenlinien schienen zu pulsieren, als würde die Route selbst atmen. Auf einem großen Bildschirm flackerte die Zoom-Konferenz -- die Gesichter von ARS, Michael Phillips und Agent Novak waren zu sehen.

Agent Novak, der I.R.A.R.A.H-Einsatzleiter, war bereits vor der Kamera, sein Blick fest und konzentriert. „Guten Morgen, alle zusammen``, begann er, seine Stimme strahlte sowohl Autorität als auch Besorgnis aus. „Unsere Mission ist klar: Wir müssen zwei Männer sicher aus der Ukraine bringen -- einen ukrainischen Pazifisten und einen russischen Deserteur. Die Theiß ist die letzte Barriere, die wir überwinden müssen, bevor wir nach Rumänien zurückkehren. Diese Region wird streng überwacht, und wir haben nur ein kleines Zeitfenster.``

Ein Raunen ging durch den Raum. Jeder wusste: Das Leben dieser Männer lag in ihren Händen.

ARS` beruhigende, fast menschlich wirkende Stimme erfüllte den Raum. „Die Fluchtroute wurde sorgfältig geplant. Wir haben die Positionen der Grenzpatrouillen erfasst und den sichersten Punkt für die Überquerung der Theiß ausgewählt. Die Kommunikation während der Mission wird über verschlüsselte Kanäle erfolgen. Michael wird via Satellitentelefon jederzeit in Kontakt bleiben.``

Michael Phillips, der als Teamleiter fungierte, nickte zustimmend. Seine Augen strahlten eine Mischung aus Zuversicht und Besorgnis aus. „Denkt daran -- das Leben dieser Männer liegt in unseren Händen. Jeder falsche Schritt könnte bedeuten, dass wir auffliegen. Also bleibt ruhig und konzentriert. Es ist wichtig, dass wir als Team zusammenarbeiten.``

Er sah in die Gesichter seiner Kollegen. Er bemerkte die Entschlossenheit -- aber auch die Unsicherheit, die wie ein Schatten über dem Raum schwebte.

„Wir haben alles vorbereitet, und ich vertraue auf eure Fähigkeiten. Wenn wir zusammenhalten, können wir diese Herausforderung meistern.``

Ein kurzes Nicken, ein ermutigendes Lächeln hier und da -- die Anspannung begann sich langsam zu lösen.

„Jetzt zu den Details``, fuhr Michael fort und wandte sich wieder der Landkarte zu. „Wir werden uns in einem kleinen Dorf am Ufer der Theiß treffen. Dort wartet der Kontaktmann auf uns, der uns zu den Männern führen wird. Die Flucht muss schnell und leise verlaufen -- keine Lichtsignale, keine Geräusche. Jeder von uns hat eine Rolle. Fragen?{\kern0pt}``

Ein paar Hände hoben sich. Michael beantwortete die Fragen mit klaren, präzisen Antworten. Als er das letzte Wort sprach, spürte er, dass das Team bereit war.

„Gut, wir haben wenig Zeit. Versammelt eure Ausrüstung und trefft euch in einer halben Stunde im Garagenbereich. Jeder weiß, was zu tun ist.``

Das Team erhob sich. Michael blieb einen Moment länger, die Karte vor sich betrachtend. Die roten Linien schienen zu pulsieren.

„ARS``, murmelte er. „Gibt es Risiken, die wir übersehen haben?{\kern0pt}``

„Alle relevanten Informationen sind in der Analyse enthalten. Die Wahrscheinlichkeiten sind günstig -- solange wir den festgelegten Zeitrahmen einhalten und keine Abweichungen vom Plan vornehmen.``

Michael atmete tief durch. Die Mission begann.

Das Team bestieg einen Flug von New York nach Bukarest. In der ersten Klasse, wo sie Platz genommen hatten, war die Atmosphäre von gespannter Erwartung durchzogen. Captain Lukas Berger, ein erfahrener Hochseekapitän mit einem tiefen Verständnis für die Gefahren der Flucht, lehnte sich leicht vor.

„Die Strömung auf der Theiß ist unberechenbar``, erklärte er mit ernstem Blick. „Wir müssen schnell und leise sein. Der Motor des Boots ist gedämpft, aber jede Bewegung kann Aufmerksamkeit erregen. Die Region wird streng überwacht.``

Julia hörte aufmerksam zu. „Jeder von uns muss seinen Teil perfekt erledigen. Unsere Fehler dürfen nicht die Freiheit anderer gefährden.``

In einer anderen Reihe unterhielten sich Michaels Doppelgänger und Dr. Neumann -- der Professor, den sie bereits in Sicherheit gebracht hatten. Ihre Unterhaltung war leise, fast gedämpft. Sie sprachen über die zunehmende Überwachung durch Technologie und die schleichende Erosion der persönlichen Freiheit.

„Transhumanismus und Technokratie gehen Hand in Hand``, murmelte Neumann. „Sie versuchen, den Diskurs durch Kontrolle zu ersticken. Technik wird dabei als Werkzeug benutzt.``

Der Doppelgänger nickte energisch. „Deshalb müssen wir zeigen, dass Freiheit etwas anderes ist als technologische Überlegenheit. Das ist die Mission von I.R.A.R.A.H. Unsere Werte sind untrennbar mit der menschlichen Würde verbunden.``

Nach der Landung in Bukarest spürten sie den Druck der bevorstehenden Aufgabe. Die Reise nach Sighetu Marmatiei, einer Grenzstadt an der Theiß, war lang -- fast neun Stunden Fahrt durch malerische, aber von den Überresten vergangener Konflikte geprägte Landschaften.

ARS übernahm die Kontrolle über die Kommunikation. „Ihr nähert euch dem Fluss``, sagte die beruhigende Stimme über das verschlüsselte Funkgerät. „Haltet euch an die geplante Route. Ich habe die Patrouillenbewegungen in Echtzeit im Blick.``

Die Landschaft war melancholisch -- weitläufige Felder, alte, verwitterte Dörfer, die Geschichten von Hoffnung und Verzweiflung flüsterten.

Am vereinbarten Ort angekommen, parkten sie das Fahrzeug an einem abgelegenen, schattigen Platz. Das dichte Unterholz schützte sie vor neugierigen Blicken.

Der leise Klang des Flusses war in der Luft zu hören. Sie bewegten sich zu Fuß durch das Dickicht. Mit jedem Schritt wuchs die Anspannung. Schließlich entdeckten sie das Boot, das von I.R.A.R.A.H bereitgestellt worden war -- versteckt zwischen hohen Bäumen, der Motor gedämpft, bereit für den Start.

Der Himmel war von bedrohlichen, dunklen Wolken bedeckt -- ein Vorbote ihrer Mission.

Sie stiegen ins Boot. Captain Berger nahm das Steuer. Die Motoren sprangen an -- ein Adrenalinstoß durchbrach die Stille.

Berger navigierte durch das kalte, dunkle Wasser der Theiß. Die Strömung war stärker, als er erwartet hatte, aber seine Erfahrung hielt ihn ruhig und konzentriert.

„Wir nähern uns dem Treffpunkt``, meldete ARS. „Die Männer sind in einem Waldstück am ukrainischen Ufer versteckt. Wir haben nur wenige Minuten, um sie zu finden.``

Alle starrten gebannt auf den Horizont. Dann -- zwei Gestalten, die aus dem Gebüsch traten. Ihre Gesichter waren gezeichnet von Erschöpfung, Angst und Hoffnung.

„Schnell, kommt an Bord!{\kern0pt}``, rief Julia und half ihnen hastig ins Boot.

Kaum waren die Männer sicher an Bord, drängte Berger mit voller Kraft zurück zum rumänischen Ufer.

Dann blitzte ein Lichtschein am Horizont auf -- eine Patrouille war in der Ferne sichtbar.

„Beeilt euch!{\kern0pt}``, rief Berger. Der Motor heulte auf, das Boot schoss durch das Wasser. Der Doppelgänger hielt einen Stoffumhang über die Seite, um sie vor den Scheinwerfern zu verbergen.

„Bleibt ruhig!{\kern0pt}``, flüsterte Julia. „Es wird alles gut. Wir sind auf dem richtigen Weg.``

Endlich erreichten sie das rumänische Ufer. Hastig zogen sie das Boot ins Dickicht. Ein lokaler I.R.A.R.A.H-Kontakt wartete bereits.

„Wir werden es versenken, damit es keine Spuren gibt``, erklärte er leise, während seine Augen unruhig die Umgebung scannten. „ARS hat uns einen sicheren Fluchtweg gezeigt. Aber wir müssen uns beeilen. Die Patrouillen könnten jede Minute hier sein.``

Sie halfen den geretteten Männern aus dem Boot und führten sie durch das Dickicht zu einem kleinen, abgelegenen Jesuitenkloster. Die Jesuiten hatten bereits Vorkehrungen getroffen -- neue Identitäten, ein sicherer Ort.

Ein älterer Priester mit ruhigem Blick und sanftem Lächeln begrüßte sie. „Ihr habt viel riskiert, um diese Männer hierher zu bringen. Sie sind nun in Sicherheit. Wir werden uns um sie kümmern.``

Julia fühlte eine Welle der Erleichterung. Der Druck auf ihren Schultern fiel ab.

Nachdem sie die Geretteten in Sicherheit gebracht hatten, machte sich das Team auf den Weg nach Ungarn. Sie wählten eine abgelegene Route über die Grenze, um nicht entdeckt zu werden.

Michael kontaktierte Julia über das Satellitentelefon. „Ihr habt es geschafft. I.R.A.R.A.H hat bestätigt, dass keine Hinweise auf eure Anwesenheit zurückgeblieben sind. Gute Arbeit.``

„Danke``, antwortete Julia. „Aber wir wissen, dass dies erst der Anfang ist.``

Sie erreichten eine kleine Stadt in Nordost-Ungarn, wo sie sich neu formieren konnten.

Julia bereitete sich darauf vor, ihre Arbeit als Sozialarbeiterin und psychologische Beraterin aufzunehmen. Martina wollte in der Archäologie Fuß fassen. Aber tief in ihnen wussten beide: Die Zeit der Ruhe würde kurz sein. Die Flucht über die Theiß war nur ein kleiner Sieg im Kampf gegen die Unterdrückung, die Wellen der Veränderung würden nicht aufhören.

In den kommenden Tagen wurden die Nachrichten über neue Repressionen in der Ukraine und Russland immer drängender. Agent Novak kontaktierte sie mit einer neuen Anweisung.

„Die Situation verschärft sich``, sagte er. „Wir haben Berichte über ein bevorstehendes groß angelegtes Militärmanöver an der Grenze. Wir müssen die Augen offenhalten und bereit sein, sofort zu handeln.``

„Das lässt uns keine Ruhe``, murmelte Julia. „Aber wir sind bereit. Immer bereit.``

Die gefährliche Mission an der Theiß war abgeschlossen. Aber die Gefahr war nicht vorbei. Ihre Reise hatte sie stärker gemacht, die Entschlossenheit in ihren Herzen brannte heller denn je.

Gemeinsam würden sie die nächste Mission annehmen, die I.R.A.R.A.H ihnen anvertraute -- egal wie herausfordernd sie auch sein mochte.

\section{Wiedersehen in Budapest}\label{wiedersehen-in-budapest}

Die kühle Morgenluft umhüllte Michael, als er vor dem Eingang des Collegium Germanicum in Rom stand. Der steinerne Bau mit seinen alten Mauern schien ihn stumm zu verabschieden. Er spürte das Gewicht der Jahre, die er hier verbracht hatte, und die Erinnerungen, die in den Wänden gespeichert waren. Die Sonne brach durch die Wolken und tauchte die Fassaden in ein sanftes, goldenes Licht. Mit einem letzten Blick auf die vertraute Umgebung ließ er die schwere Holztür hinter sich zufallen.

Draußen wartete Maria auf ihn, in einen leichten Mantel gehüllt, der sanft im Morgenwind flatterte. Ihr Blick war ruhig, doch in ihren Augen lag ein Ausdruck, der Michael für einen Moment innehalten ließ. Es war der Blick einer Frau, die viel zu sagen hatte -- aber die Worte in der Stille gefangen hielt.

„Also, es ist soweit``, sagte sie, ihre Stimme weich, aber durchzogen von einer Melancholie, die die Luft um sie herum schwer machte. „Du reist ab und lässt alles hinter dir.``

Michael nickte langsam. „Ja, es ist Zeit. Es gibt viel zu tun in Budapest.`` Er hielt kurz inne und sah ihr tief in die Augen. „Ich hoffe, du weißt, dass ich immer noch an dich denke -- und an alles, was wir geteilt haben.``

Ein leichtes Lächeln huschte über Marias Lippen, doch ihre Augen verrieten eine tiefere Geschichte. „Pass auf dich auf, Michael. Es ist gut zu wissen, dass du die Familie nicht vergisst.``

Für einen Augenblick schien es, als läge in ihren Worten mehr, als sie aussprach. Michael wusste, dass die Andeutung über die Vergangenheit und die Identität des Doppelgängers wie ein unausgesprochener Schatten zwischen ihnen stand. Er neigte den Kopf, und ohne ein weiteres Wort wandte er sich ab, um zum Taxi zu gehen, das ihn zum Flughafen bringen würde.

-\/-\/-

Der Flug nach Budapest verlief ruhig. Während er durch das Fenster auf die unter ihm vorbeiziehende Landschaft blickte -- die Alpen, dann die weite ungarische Tiefebene -- kreisten seine Gedanken um das, was vor ihm lag und was er in Rom zurückgelassen hatte. Der Gedanke an Maria, an die unausgesprochenen Worte, die zwischen ihnen schwebten, drängte sich in den Vordergrund.

Als das Flugzeug auf dem Flughafen Budapest landete, durchzog ein vertrautes Kribbeln der Vorfreude seinen Körper. Es war nicht nur eine neue Aufgabe, die auf ihn wartete -- sondern auch die Möglichkeit, endlich Klarheit zu schaffen.

Ein schwarzer Wagen wartete auf ihn. Die Fahrt durch die Stadt führte ihn vorbei an der Donau und den prächtigen Gebäuden, die im goldenen Abendlicht erstrahlten. Das Wasser funkelte in der Dämmerung und schien die Vergangenheit und Zukunft zu reflektieren. Er war auf dem Weg zu der kleinen Wohnung, die Julia und Martina mittlerweile ihr Zuhause nannten.

Als der Wagen vor dem alten Gebäude hielt, atmete Michael tief durch. Alte Bäume umrahmten den Eingang, die vertrauten Klänge der Stadt erfüllten die Luft. Er drückte die Klingel und wartete.

-\/-\/-

Die Tür öffnete sich, und Martina begrüßte ihn mit einem warmen Lächeln, das ein Stück der Anspannung von seinen Schultern nahm. „Michael, schön, dich zu sehen. Komm rein, wir haben auf dich gewartet.``

Im Wohnzimmer herrschte eine heimelige Atmosphäre. Der Duft von frisch gebrühtem Kaffee erfüllte den Raum, auf dem Couchtisch standen Kuchen und Gebäck bereit, die wie kleine Kunstwerke aussahen. Julia kam aus der Küche, ein Tablett in den Händen, das sie mit einem strahlenden Lächeln absetzte.

„Endlich bist du da``, sagte sie und sah Michael mit einem Ausdruck der Erleichterung an. „Setz dich, nimm dir einen Kaffee. Wir haben viel zu besprechen.``

Michael nahm Platz auf einem der alten Sofas, die mit einem liebevollen, aber abgewetzten Stoff bezogen waren. Die Gemütlichkeit des Raumes gab ihm ein Gefühl von Heimat, das er lange vermisst hatte.

Kurz darauf betrat auch der Doppelgänger den Raum. Er wirkte entspannt -- aber auch ein wenig nervös, als er Michael gegenüber Platz nahm.

„Willkommen in Budapest``, sagte er mit einem leichten Lächeln, das sowohl Offenheit als auch Unsicherheit verriet. „Es ist schön, die ganze ‚Familie` beisammen zu haben.``

Das Wort „Familie`` klang für Michael ungewohnt vertraut. Der Ausdruck im Gesicht des Doppelgängers schien für einen kurzen Moment eine tiefere Verbundenheit zu verraten. Michael erwiderte das Lächeln -- während ihm ein Gedanke durch den Kopf schoss: *War es wirklich möglich, dass dieser junge Mann sein Sohn war?*

„Danke``, sagte Michael und nahm einen Schluck von seinem Kaffee, der warm und aromatisch war. „Es fühlt sich gut an, hier zu sein.``

-\/-\/-

Sie plauderten über Belangloses -- die Stadt, die Arbeit, das Leben in Budapest. Doch zwischen den Gesprächen schwebte eine ungesprochene Frage im Raum, eine Spannung, die weder Julia noch Martina zu lösen schienen. Der Doppelgänger warf Michael hin und wieder Blicke zu, die von einem intensiven Interesse durchdrungen waren, als ob er eine Bestätigung suchte, die Michael noch nicht ausgesprochen hatte.

Das Abendessen verlief entspannt. Die Gespräche waren leicht und voller Lachen. Doch als sie sich schließlich in der Dämmerung auf dem Balkon niederließen, lag über der Szene eine Art von stiller Verständigung. Michael wusste, dass es an der Zeit war, Antworten zu suchen -- und dass er sie möglicherweise bereits vor sich hatte.

„Manchmal``, begann er leise, als die ersten Sterne am Himmel aufleuchteten, „führt uns das Leben auf unerwartete Wege, die wir erst später verstehen.`` Er sah den Doppelgänger an und bemerkte, dass dieser seine Worte aufmerksam verfolgte. „Und manchmal treffen wir Menschen, die uns zeigen, dass es mehr Verbindungen gibt, als wir zunächst glauben.``

Der Doppelgänger sagte nichts. Aber er nickte.

Die Nacht senkte sich über Budapest, und die Lichter der Stadt blinkten wie kleine Funken, die Erinnerungen in die Dunkelheit warfen. Die Blicke, die sie tauschten, sprachen Bände, während die Stille des Abends sie umgab. Michael wusste, dass es Zeit brauchen würde, um die Wahrheit vollständig auszusprechen -- aber in diesem Moment genügte es, dass sie zusammen waren.

Die Familie hatte eine neue Dimension bekommen -- eine, die er nicht erwartet, aber vielleicht immer erhofft hatte.

Das sanfte Rauschen der Donau war in der Ferne zu hören. Michael wusste: Dies war erst der Anfang einer langen und spannenden Reise.

\end{document}
